%% Livret de corrigés d\'exercices — généré pour rendu agrégé.
%% Niveau : N11
%% Séries incluses : R, AE, F, A, E

\documentclass[a4paper,11pt,twoside]{article}

%% ── Paquets externes (noyau pour atome) ───────────────
\usepackage[utf8]{inputenc}
\usepackage[T1]{fontenc}
\usepackage{lmodern}
\usepackage[french]{babel}
\usepackage{geometry}
\usepackage{amsmath}
\usepackage{amssymb}
\usepackage{etoolbox}
\usepackage{xkeyval}
\usepackage{xifthen}
\usepackage{environ}
\usepackage{forloop}
\usepackage{xcolor}
\usepackage{ninecolors}
\usepackage{multirow}
\usepackage{colortbl}
\usepackage{tabularray}
\usepackage{multicol}
\usepackage[inline]{enumitem}
\usepackage{graphicx}
\usepackage{adjustbox}
\usepackage{tikz}
\usepackage[theorems,breakable,skins]{tcolorbox}
\usepackage{pifont}

%% ── Marges (alignées sur les livrets) ───────────────
\newgeometry{left=35pt,right=35pt,top=65pt,bottom=70pt}

%% ── Initialisations de bibliothèques ─────────────────
\UseTblrLibrary{booktabs}
\UseTblrLibrary{varwidth}

%% ── Définitions du paquet seqenseigne (inlinées) ─────
%%    Fermeture transitive des macros effectivement utilisées
%%    par l'atome — au lieu de \usepackage{seqenseigne}
%%    qui chargerait ~110 paquets.
\makeatletter
\makeatother

%% ── Macros transverses ───────────────────────────────
\providecommand{\acompleter}[1]{\;\underline{\phantom{#1}}\;}

\usepackage{titlesec}
\usepackage{fancyhdr}
\usepackage{lastpage}

\geometry{vmargin=70pt,hmargin=50pt,headheight=50pt,headsep=15pt,footskip=20pt}

\pagestyle{fancy}
\fancyhf{}
\lhead{\leftmark}
\rhead{\rightmark}
\rfoot{Page \thepage/\pageref{LastPage}}
\lfoot{Le \today}
\renewcommand{\sectionmark}[1]{\markboth{#1}{}}
\renewcommand{\subsectionmark}[1]{\markright{#1}}
\renewcommand{\headrulewidth}{0.4pt}
\renewcommand{\footrulewidth}{0pt}

\titleformat{\section}[frame]{\normalfont}{\filright\footnotesize\enspace SECTION \thesection\enspace}{30pt}{\huge\bfseries\filcenter}
\titleformat{\subsection}[block]{\normalfont\sffamily}{\thesubsection}{.5em}{\titlerule[2pt]\\[1em]\bfseries\sc\Large\filcenter}[\vspace{2ex}\titlerule]
\titleformat{\subsubsection}{\normalfont\bfseries}{\thesubsubsection}{1em}{}

\setcounter{tocdepth}{2}

\begin{document}

\providecommand{\theHExoNum}{\theExoNum}

\seqCreeCompteurs
\seqSetCodeNiveau{N11}

\thispagestyle{empty}
\begin{center}
\vspace*{2em}
{\Huge\bfseries Livret de corrigés d\'exercices}\\[2em]
{\Large Collège — Document enseignant}\\[1em]
{\LARGE Classe de 4\ieme{}}\\[1em]
{\large \emph{Séries : R, AE, F, A, E}}\\[3em]
\end{center}

\begin{itemize}
  \item Représentations d'un nombre
\end{itemize}

\newpage

\thispagestyle{empty}
\tableofcontents
\newpage

\cleardoublepage
\seqSetColorsTheme{nombres}
\section{Nombres}

\seqSetCodeSequence{S01}
\subsection{Représentations d'un nombre}

\subsubsection{Révisions / activités d'entrée}

\begin{description}
  \item[Exercice 1.]\quad
Corrigé EA1

\end{description}

\subsubsection{Série F (formation)}

\begin{description}
  \item[Exercice 1.]\quad
Corrigé F1

\end{description}

\subsubsection{Série A}

\begin{description}
  \item[Exercice 1.]\quad
Corrigé A1

\end{description}

\subsubsection{Série E (entraînement)}

\begin{description}
  \item[Exercice 1.]\quad
Corrigé E1

\end{description}

\end{document}